%
% pdfTeX 3.141592653-2.6-1.40.26 (TeX Live 2025/dev/Debian)
% 0001multiples_of_3_or_5.tex
%
% Copyright 2026 @GitHub JoergEm
%
% This program is WITHOUT ANY WARRANTY; without even the implied
% warranty of MERCHANTABILITY or FITNESS FOR A PARTICULAR PURPOSE.
%

\documentclass{article}
\usepackage{siunitx}
\ExplSyntaxOn

\NewDocumentCommand{\setvariable}{mmm}
 {
  \prop_gclear_new:c { g_variabeln_var_#1_prop }
  \prop_gput:cnn { g_variabeln_var_#1_prop } { value } { #2 }
  \prop_gput:cnn { g_variabeln_var_#1_prop } { unit } { #3 }
 }

\NewDocumentCommand{\printvariable}{sm}
 {
  \IfBooleanTF{#1}
   { \num { \__variabeln_get:nn { #2 } { value } } }
   { \exp_args:Nnx \SI { \__variabeln_get:nn { #2 } { value } }
       { \__variabeln_get:nn { #2 } { unit } } }
 }

\DeclareExpandableDocumentCommand{\getvalue}{m}
 { \__variabeln_get:nn { #1 } { value } }

\cs_new:Npn \__variabeln_get:nn #1 #2
 { \prop_item:cn { g_variabeln_var_#1_prop } { #2 } }

\cs_set_eq:NN \fpeval \fp_eval:n

\cs_new:Npn \__euler_sum:nn #1 #2
 {
  \fp_eval:n
   {
    floor( (#1 - 1) / #2 ) * ( floor( (#1 - 1) / #2 ) + 1 ) / 2 * #2
   }
 }

\NewDocumentCommand{\multiplesofthreeorfive}{mm}
 { \__euler_sum:nn { #1 } { #2 } }

\ExplSyntaxOff

\setvariable{nummer}{1000}{}

\begin{document}

\section*{Project Euler -- Problem 1}

Gesucht ist die Summe aller Vielfachen von~3 oder~5 unterhalb von
\printvariable*{nummer}.

\bigskip
\noindent
Die Gaußsche Summenformel liefert die Summe aller Vielfachen von $k$
unterhalb von $N$:
\[
  S(k,N)
  = \left\lfloor\frac{N-1}{k}\right\rfloor
    \cdot \frac{\left\lfloor\frac{N-1}{k}\right\rfloor + 1}{2}
    \cdot k
\]

\bigskip
\noindent
\begin{tabular}{@{}ll@{}}
  \textbf{Vielfache von 3:}                    & $\multiplesofthreeorfive{\getvalue{nummer}}{3}$  \\[4pt]
  \textbf{Vielfache von 5:}                    & $\multiplesofthreeorfive{\getvalue{nummer}}{5}$  \\[4pt]
  \textbf{Vielfache von 15 (Doppelzählungen):} & $\multiplesofthreeorfive{\getvalue{nummer}}{15}$ \\[4pt]
  \textbf{Summe:}                              & $\fpeval{
    \multiplesofthreeorfive{\getvalue{nummer}}{3}
    + \multiplesofthreeorfive{\getvalue{nummer}}{5}
    - \multiplesofthreeorfive{\getvalue{nummer}}{15}
  }$                                                                                               \\
\end{tabular}

\end{document}